%! TeX program = lualatex
\documentclass{book} 

\title{Algebry Toeplitza i twierdzenie o indeksie}
\author{Weronika Jakimowicz}
\date{20.03.2026}

\usepackage{../template}

\begin{document}
% \maketitle 

Kto lubi kompleksy łańcuchowe?

Powiemy, że ograniczony operator $F:V\to W$ między dwoma przestrzeniami Banacha jest Fredholmowski, jeśli jego jądro i kojądro mają skończony wymiar. Ciąg
$$\begin{tikzcd}
  K:&0\arrow[r] & V\arrow[r, "F" above] & W\arrow[r] & 0
\end{tikzcd}$$
jest kompleksem łańcuchowym - obraz strzałki jest zawarty w jądrze następnej strzałki. Możemy więc liczyć jego kohomologie! $H^0(K)=\ker F$, $H^1(K)=\coker F$. Indeks $F$ to będzie charakterystyka Eulera tego kompleksu, czyli alternująca suma wymiarów jego kohomologii:
$$ind(F)=\chi(K)=\dim(H^0(K))-\dim(H^1(K)).$$

Tak się składa, że czasami operatory Toeplitza są Fredholmowskie, ale wpierw potrzebujemy powiedzieć co to wogóle jest algebra i operator Toeplitza.

$\lambda$ - miara Haara na $S^1$

Dla $n\in\Z$ definiujemy operator
$$\epsilon_n: S^1\to S^1$$
$$\epsilon_n(z)=z^n$$
Interesują nas dwie podalgebry (jeszcze nie powiem w czym)
$$\Gamma=\span\{\epsilon_n\;:\;n\in\Z\}$$
$$\Gamma_+=\span\{\epsilon_n\;:\;n\in\N\}$$

Kilka faktów o $\Gamma$:
\begin{itemize}
  \item $\epsilon_n$ tworzą ortonormalną bazę $L^2S^2$, bo to proste wyliczenie
    $$\int\epsilon_n\overline{\epsilon_m}\lambda= \int_0^{2\pi} e^{itn}\overline{e^{itm}}dt=\int_0^{2\pi}e^{it(n-m)}=\begin{cases}1 & n=m\\0&n\neq m\end{cases}$$
  \item Stone-Weierstrass mówi, że jest gęsta w $CS^1$, a ponieważ $CS^1$ jest gęste w $L^pS^1$ to $\Gamma$ też (dla nas ważne, że w $L^2S^1$, bo mamy bazę Hilberta)
  \item jest *-podalgebrą $CS^1$ (addytywne bo span, mnożenie bo $z^nz^m=z^{n+m}$, inwolucja to minus w wykładniku, C* to S-W punkt wyżej) [{\color{red}ĆWICZENIE?}]
\end{itemize}

\begin{definition}{współczynniki Fouriera}{}
  Dla $f\in L^pS^1$ definiujemy $n$-ty współczynnik Foueriera jako
  $$\hat{f}(n)=\int f(x)\overline{x^n}d\lambda(x)$$
\end{definition}

\begin{example}
  $f(x)=x^m$, wtedy $\hat{f}(n)=\begin{cases}1 & n=m\\0 & wpp\end{cases}$ bo to jest dokładnie to samo co liczyliśmy wcześniej 
\end{example}

\begin{definition}{Hardy space}{}
  $$H^p(S^1)=\{f\in L^p(S^1)\;:\;\hat{f}(n)=0\text{ dla }n<0\}$$
\end{definition}

Nas będzie interesować $H^2(S^1)$.

Fun fact: $\Gamma$ jest gęste i ortonormalne w $L^2S^1$, a $\Gamma_+$ - w $H^2S^1$.

Mamy rzut ortogonalny
$$P:L^2S^1\to H^2S^1.$$
Dla $\phi\in L^\infty S^1$ definiujemy operator mnożenia jako
$$M_\phi(f)=\phi f$$
a \buff{operator Toeplitza} jako rzut operatora mnożenia
$$T_\phi=PM_\phi.$$
Funkcję $\phi$ często nazywa się \buff{symbolem} operatora $T_\phi$.

\begin{lemma}{}{}
  $T_\phi$ is compact iff $\phi=0$
\end{lemma}

\begin{proof}\color{red} 3.5.8 Murphy

  $v:=M_{\epsilon_1}$ (przesuwa bazę), $u:=v\restriction{H^2S^1}$.
\end{proof}

\begin{lemma}{}{}
  $\phi\in CS^1$, $\psi\in L^\infty S^1$ wtedy $T_\phi T_\psi-T_{\phi\psi}$ and $T_\psi T_\phi-T_{\psi\phi}$ są zwarte
\end{lemma}

\begin{theorem}{Toeplitza o indeksie}{}
  Jeśli $\phi$ jest ciągłe, to $T_\phi$ jest Fredholmowskie oraz 
  $$ind(T_\phi)=-winding(\phi),$$
  gdzie $winding(\phi)$ to ilość nawinięć $\phi$ na $S^1$.
\end{theorem}

% Zanim to pokażemy potrzebujemy szybko zanurzyć stópki w świecie operatorów zwartych.

\begin{proof}\color{red}
  charakteryzacja Atkinsona ble ble ble 3.5.12
\end{proof}

Przypomnienie:
\begin{description}
  \item[spektrum algebry] (aka spektrum Gelfanda) $\Spec(A)=\{\text{niezerowe *-homomorfizmy }\phi:A\to \C\}$ (akurat zgadza się ze spektrum maksymalnym $A$)
  \item[widmo elementu] $\sigma(a)=\{\lambda\in\C\;:\;a-\lambda\text{ nieodwracalne}\}$
  \item[widmo operatora] $\sigma(T)=\{\lambda\;:\;T-\lambda\text{ nieodwracalne}\}$
\end{description}

$B(H)$ niech będą operatory ograniczone, $K(H)$ - operatory zwarte, czyli domknięcie obrazu kuli jednostkowej przez nie jest zwarte. 
$$\pi:B(H)\to B(H)/K(H)$$

\begin{definition}{widmo właściwe operatora (ang. essential spectrum)}{}
  Niech $T$ będzie operatorem ograniczonym na przestrzeni Hilberta $H$. 
  Widmo zwarte operatora $T$ to spektrum obrazu rzutu $T$ na $B(H)/K(H)$.
\end{definition}

Intuicja: widmo operatora mówi kiedy nie da się odwrócić $T-\lambda$, widmo właściwe powie kiedy nawet dodanie jakiegoś małego (czyli zwartego) operatroa nie pomoże.

% \begin{theorem}{}{}
%   $T_1$, $T_2$ - samosprzężone operatory na ośrodkowej przestrzeni Hilberta (którą jest np. $L^pS^1$, która nas będzie potem interesować) mające to samo widmo właściwe. Wtedy istnieje unitarny automorfizm $U$ dla którego $T_1-U*T_2U$ jest zwartym operatorem.
% \end{theorem}
%
% Dowód używa maszynki, którą Weyl i von Neumann dawno temu wytworzyli:
% \begin{theorem}{Weyl-von Neumann}{}
%   Każdy ograniczony, samosprzężony operator na ośrodkowej przestrzeni Hilberta jest dowolnie małym zwartym potrząśnięciem operatora diagonalnego (macierz jest diagonalna).
% \end{theorem}
%
% \begin{proof}
%   Tego twierdzenia wyżej.
%
%   $T_1$ i $T_2$ to potrząsnięcia diagonalnego operatora zwartego mające to samo spektrum.
% \end{proof}





\end{document}
